\scanpage{51}
9. \uline{Cas $S$ de car} 2 (i.e. $2\,1_S=0$)

Dans ce cas \deleted{$2F\simeq\check{\mathbb V}(L)$}, on a iso can.
\[2F\simeq\check{\mathbb V}(L)\tag*{(95)}\]
ou ce qui revient au $=$, la condition $\mathrm{tr}(e_B)=2$ sur une forme linéaire $\mathrm{tr}$ sur $B,\unsure{e_B}$ muni de $e_B$ \deleted{implique \unsure{$\mathrm{tr}$}} \inserted{signifie}
\[\mathrm{tr}(e_B)=0\tag*{(96)}\]
de sorte que la forme trace $\mathrm{tr}$ (\deleted{\unsure{ou}} ou la section $\underline b$.) est une section de $L$
\[\underline b=\mathrm{tr}\in\Gamma(S,L)\tag*{(97)}\]
canoniquement associée à \deleted{\unsure{tte|la}} structure d'alg. quadratique sur $B$. \deleted{Ceci dit, on a ceci} \inserted{D'ailleurs,} le discriminant est donné par
\[D=\underline b^2\tag*{(98)}\]
car en termes d'un scindage \inserted{$s$} on aura \deleted{\unsure{$b(s)$}}
\[\deleted{\unsure{b_s}}\,b(s)=\underline b\qquad\text{(indépendant de $s$ !)}\tag*{(99)}\]
et
\[D=b^2-4c=b^2=\underline b^2\qquad\text{(car $4=0$)}\]

\uline{NB} Ici l'équation $x^2=D$ s'écrit donc $(x^{\deleted{\unsure{2}}}-\underline b)^2=0$, et donne un sous-schéma de $\check{\mathbb V}(L)$ isom. à $\alpha_2(L)$, et sans grande relation avec $\Delta$\ldots

En concl, \inserted{en plus de $\underline b\in\Gamma(S,L)$,} le torseur $F$ sous $L$, et sa classe
\[f\in H^1(S,L)\tag*{(100)}\]
est un \unsure{nouvel} invariant de la situation. L'involution
% archive sheet number: 50
